% Nota introductoria de esta versión (no es texto del documento original).
\section{Sobre esta versión}\label{sec:intro}

Este documento reúne únicamente los Derechos Básicos de Aprendizaje (DBA, versión 1, 2016)
de \textbf{física} de Ciencias Naturales para los grados 9, 10 y 11:

\begin{itemize}[nosep]
  \item \textbf{Grado 9}, DBA 1: descripción y predicción del movimiento en un marco de referencia inercial.
  \item \textbf{Grado 10}, DBA 1: fuerzas, equilibrio y leyes de Newton.
  \item \textbf{Grado 10}, DBA 2: conservación de la energía mecánica.
  \item \textbf{Grado 11}, DBA 1: sonido y luz como fenómenos ondulatorios.
  \item \textbf{Grado 11}, DBA 2: fuerzas eléctricas y magnéticas.
  \item \textbf{Grado 11}, DBA 3: corriente y voltaje en circuitos resistivos.
\end{itemize}

Se conserva la numeración original de cada DBA dentro de su grado; los números que faltan
corresponden a DBA de química o biología, que no se incluyen. El texto completo de todos los
grados (1 a 11) y áreas está disponible como referencia en \texttt{raw/grado\_NN.txt}.

Cada DBA tiene tres partes: el \emph{enunciado} (el aprendizaje estructurante), las
\emph{evidencias de aprendizaje} (indicios que muestran si se está alcanzando) y el
\emph{ejemplo} (concreta las evidencias). Las figuras aparecen en su lugar, como en el
original, con una descripción como pie; son archivos PNG en \texttt{figs/}. Varios ejemplos
consisten solo en una imagen: su pie describe la situación física que plantean.
